%%%%%%%%%%%%%%%
%% cubic-code.tex
%%%%%%%%%%%%%%%

The toric code~\cite{Kitaev2003Fault-tolerant} 
is usually defined on a planar graph with qubits residing on edges.
The star operator and plaquette operators 
are defined according to the data of the graph.
The model is thus a priori lattice dependent.
However, a certain set of important properties of the model
turns out to be lattice independent.
Particularly the ground-state subspace 
conveys a structure that depends only on the topology of the underlying space, 
insensitive to the microscopic detail.
If the underlying space is a genus $g$ (oriented) surface,
the degeneracy is $4^g$.
Even if the Hamiltonian is perturbed,
the energy splitting of lowest $4^g$ states
is exponentially small in the system size
as long as the perturbation is small enough~%
\cite{WenNiu1990GroundState,
Kitaev2003Fault-tolerant,
BravyiHastingsMichalakis2010stability}.
It is a signature of \emph{topological order}.


Authors have used the term `topological quantum order' to mean
all or a part of the following properties:
Ground-state degeneracy as a function of topology,
no spontaneous symmetry breaking,
anyonic particle content~%
\cite{
WenWilczekZee1989Chiral,
Wen1989Degeneracy,
Einarsson1990Fractional,
WenNiu1990GroundState,
ReadSachdev1991LargeN,
Wen1991SpinLiquid},
robust edge modes against 
perturbations~\cite{KaneMele2005,FuKaneMele2007Topological,HasanKane2010TIReview},
locally indistinguishable ground states~\cite{
BravyiHastingsMichalakis2010stability,
BravyiHastings2011short,
MichalakisPytel2011stability},
and topological entanglement entropy~\cite{KitaevPreskill2006Topological,LevinWen2006Detecting}.
Here, we take the local indistinguishability of ground states
as a definition of topological quantum order.
The local indistinguishability is precisely the one used in Chapter~\ref{chap:alg-theory},
as well as in the proof of gap stability 
results~\cite{BravyiHastings2011short,
MichalakisPytel2011stability}.
See Lemma~\ref{lem:local-tqo=exact}.
We ask what quantum phases are possible in the class of code Hamiltonians.


An important example of topological quantum order 
presented in this thesis is \emph{cubic code}.
In this chapter we explain how the model is found,
and study its consequences.
The cubic code is an exact local additive code with translation symmetry 
on the simple cubic lattice, 
where the exactness is as in Chapter~\ref{chap:alg-theory}.
This model is a gapped unfrustrated spin Hamiltonian,
and is topologically ordered in the very sense we just defined,
but breaks many aspects of conventional models of topological order.
Most prominently, the excitations or charges of the cubic code
\emph{cannot} be interpreted as particles
since they are immobile;
the hopping term appears only after $L$-th
or higher order perturbation theory, where $L$ is the linear system size.
The immobility implies that the system spends
quite a long time to reach its thermal equilibrium
in response to environment's change.
(Later, we will give quantitative statements regarding this.)
Moreover, due to the topological quantum order,
the phase is stable with respect to arbitrary but small perturbations;
the immobility of excitations is also protected~\cite{BravyiHastingsMichalakis2010stability}.


The immobility of the charges translates 
into the theory of quantum error correcting codes
as the absence of string logical operators.
We call this property by \emph{no-strings rule}.
The notion of string logical operators might be intuitive
if one imagines the toric codes in two or three 
dimensions~\cite{DennisKitaevLandahlEtAl2002Topological}.
However, this intuition is too model specific;
the strings in discrete lattices are not well-defined objects
since a set of points in the lattice does not in general have
a well-defined dimensionality.
We overcome this issue by defining \emph{string segments}
that capture characteristics that are responsible for the mobility of charges.

We proceed by translating conditions for the absence of the logical string segments
into our algebraic framework,
in order to systematically search for codes without logical strings;
the cubic code is not an ad hoc model.
By Corollary~\ref{cor:annT-3d-characteristic-dimension-0}%
~\cite{Yoshida2011feasibility},
the characteristic dimension of the cubic code must be 1
and the degeneracy must generally grow with the system size.
An explicit formula for the degeneracy is given.
Additionally, a real-space renormalization computation is presented.
It seems that the cubic code is a fixed point of a certain unconventional kind.
Namely, the model is a direct sum of two daughter models $A$ and $B$ at a coarse-grained lattice
where one daughter model $A$ is the same as the original model, but the other $B$ is not.
The model $B$ produces two copies of itself at a further coarse-grained lattice.
The renormalization group flow continues to branch.
Next, we compute the thermal partition function,
and show that there is no finite temperature phase transition.

\section{String segments and no-strings rule}

The most important property of one-dimensional objects
is that a finite part of it has two disjoint boundary points.
An intuitive role of string operator 
is to move an excitation at its one boundary point
to another boundary point.
Essentially, it is a concatenation or juxtaposition of hopping operators.
%
The hopping operator is, as a whole, a finitely supported operator.
Therefore, if the hopping operator acts on the vacuum (ground state),
the overall effect is to create a trivial or neutral charge,
which consists of two spatially separated charges.
When either of the two charges is neutral by itself,
the hopping is meaningless 
because trivial charges can be annihilated locally
and created anywhere arbitrarily;
only the hopping of nontrivial charges is important.
Now we can define string segments and their absence.
\begin{defn}
\cite{Haah2011Local,BravyiHaah2011Energy}
A {\bf string segment} is a finitely supported Pauli operator
that creates excitations contained in the union of two finite boxes
of {\bf width} $w$.
The string segment is {\bf nontrivial} 
if the charge contained in one of the boxes is nontrivial.
The distance between the boxes is the {\bf length} of the string segment.
We say a model obeys {\bf no-strings rule}
if the length of any nontrivial string segment of width $w$
is bounded by $\alpha w$ for some constant $\alpha \ge 1$.
\label{defn:no-strings}
\end{defn}
\noindent
The no-strings rule may seem too strong than necessary;
why do we need an upper bound by a \emph{linear} function?
It is rather a technicality 
that is necessary to prove a logarithmic energy barrier theorem
in Section~\ref{sec:log-barrier}.
However, our definition seems sharp yet broad enough to derive further results.


An immediate consequence of the no-strings rule
in the case of a translation-invariant code Hamiltonian
is that there are infinitely many charges.
If there were only finitely many,
then in the sequence of all translations of a charge $c$
there would be an equivalent charge $c'$.
It means $c - c'$ is neutral 
and therefore $c-c'$ is created by a finitely supported operator,
which is a nontrivial string segment.
This string segment can be juxtaposed many times
to give arbitrarily long string segments with fixed width.
This is a contradiction to the no-strings rule.
Also, in view of Corollary~\ref{cor:annT-3d-characteristic-dimension-0}
and Lemma~\ref{lem:k-growing-sequence-L},
a translationally invariant three-dimensional exact code Hamiltonian that obeys the no-strings rule,
must have a growing ground-state degeneracy with respect to the system size
when defined on lattices with periodic boundary conditions.


\section{Search for models}

We wish to find an interesting class of models
that are relatively handy to deal with.
The following lemma shows such a family.
Recall that for an exact code Hamiltonian,
we have $\ker \sigma^\dagger \lambda = \im \sigma$.
\begin{lem}
The generating matrix of a code Hamiltonian
\[
\sigma =
\begin{pmatrix}
f & 0 \\
g & 0 \\
0 & g' \\
0 & f' \\
\end{pmatrix},
\]
where $f,g,f',g' \in R$, elements of group algebra of $\ZZ^D$,
is \emph{exact}
if and only if 
\[
\gcd(f,g)=1,
\quad \quad
g' = \bar g,
\quad
 f' = -\bar f 
\]
up to units of $R$.
\label{lem:two-terms-CSS-exactness-condition}
\end{lem}
\noindent
It is a generalization of $\sigma$ in Remark~\ref{rem:long-string}.
The 2D toric code of Example~\ref{eg:2d-toric} falls into this form.
\begin{proof}
Observe that there is no mixing between 
the $\sigma^x$ part (the upper two rows of $\sigma$)
and the $\sigma^z$ part (the lower two rows of $\sigma$).
So $\ker \sigma^\dagger \lambda = \im \sigma$ can be checked
separately.
That is, we need to verify whether the kernel of
\[
\epsilon_z = \begin{pmatrix} \bar f & \bar g \end{pmatrix}
\]
is generated by $\begin{pmatrix} g' \\ f' \end{pmatrix}$,
and whether the kernel of
\[
\epsilon_x = \begin{pmatrix} \bar g' & \bar f' \end{pmatrix}
\]
is generated by $\begin{pmatrix} f \\ g \end{pmatrix}$.
If $f = f_1 h$ and $g = g_1 h$ where $\gcd(f,g)=h$,
then $\begin{pmatrix} \bar g_1 \\ -\bar f_1 \end{pmatrix} \in \ker \epsilon_z$.
We have $\begin{pmatrix} \bar g_1 & -\bar f_1 \end{pmatrix}
 = r \begin{pmatrix} g' & f' \end{pmatrix}$ for some $r \in R$.
Since $\gcd(g_1, f_1)=1$, $r$ must be a unit.
Hence, we may set $g' = \bar g_1$ and $f' = -\bar f_1$.
Then, $\begin{pmatrix} f_1 \\ g_1 \end{pmatrix} \in \ker \epsilon_x$,
so $h$ is a unit. The converse is straightforward.
\end{proof}

We stay with the form of $\sigma$ as in
Lemma~\ref{lem:two-terms-CSS-exactness-condition}
for its simplicity.
Now we consider the no-strings rule.
The equivalence classes of topologically nontrivial charges
is $\coker \epsilon$ by Theorem~\ref{thm:charge-equals-torsion},
which decomposes as $\coker \epsilon_z \oplus \coker \epsilon_x$.
Since the two summands are related by the antipode map,
it suffices to consider $\coker \epsilon_x \cong R/(f,g)$ only.
A necessary condition for the no-strings rule is that
any fractal generator should \emph{not} consist of two terms.
For Theorem~\ref{thm:structure-2d-ann-coker-epsilon},
the minimal dimension we should be interested in is 3;
$ R = \FF_2[x^{\pm 1},y^{\pm 1},z^{\pm 1}]$.
Let $I = (f,g)$, the ideal of $R$ generated by $f$ and $g$.
If, for example, $f$ factorizes as $f=(1+x)f_1$ and $f_1 \notin I$,
then $(1+x)$ is a string generator.
So it is necessary that the Gr\"obner basis of the ideal $(f,g)$ does not have
any two-term factor such as $x+1$, $y+1$, or $z+1$.
In order to have a degenerate code Hamiltonian 
(Corollary~\ref{cor:unit-characteristic-ideal-means-nondegeneracy}),
we have to have a non-unit associated ideal.
For simplicity, we demand that $x=y=z=1$ is a root of $f$ and $g$
so that the associated ideal is contained in a maximal ideal $(x+1,y+1,z+1)$.

We further assume that $f$ and $g$ have exponents $0$ or $1$ in each variable.
Thus, $f$ and $g$ are linear combinations of $1,x,y,z,xy,yz,zx,xyz$ over $\FF_2$.
Naively there are $2^{16}$ possibilities, but they are not all different.
For example, three choices $(f(x,y,z),g(x,y,z))$, $(x f(\frac{1}{x}),y,z), x g(\frac{1}{x},y,z))$,
and $(f(y,z,x),g(y,z,x))$ define the same models 
because they are related by the reflection about $yz$-plane
or the $\pi/3$ rotation about $(1,1,1)$-axis.
Up to these symmetries of the unit cube, there are 392 pairs of $f,g$.
An exhaustive search gives
10 models in Table~\ref{tb:10-models-wo-string} that satisfy all our requirements.
The most symmetric model defined by
\[
f = 1+ x + y + z, \quad \quad g = 1+xy+yz+zx
\]
will be called {\bf cubic code}.

\begin{table}[htbp]
\centering
\begin{tabular}{c|c|c}
\hline
   & $f$ & $g$ \\
\hline
1 & $1+x+y+z$ & $1+x y+x z+y z$ \\
2 & $x+y+z+y z$ & $1+y+x y+z+x z+x y z$ \\
3 & $1+x+y+z$ & $1+x z+y z+x y z$ \\
4 & $1+x+z+y z$ & $1+y+x y+x z$ \\
5 & $1+x+z+y z$ & $y+z+x z+y z$ \\
6 & $1+x+y+z$ & $1+y+x z+y z$ \\
7 & $1+x+y+z$ & $1+z+y z+x y z$ \\
8 & $1+x+z+y z$ & $1+y+x y+z+x z+y z$ \\
9 & $1+x+y+z$ & $x y+z+x z+y z$ \\
10 & $1+x+z+y z$ & $1+y+x y+x z+y z+x y z$ \\
\hline
\end{tabular}
\caption{Complete list of cubic codes. 
Each pair of polynomials defines an exact degenerate code Hamiltonian
according to Lemma~\ref{lem:two-terms-CSS-exactness-condition}.
They potentially obey the no-strings rule.
Especially, code 1 indeed obeys the no-strings rule as proven in Section~\ref{sec:cubic-code}.
The table is exhaustive, up to symmetries of simple cubic lattice and local symplectic transformations,
under the following criteria:
(i) $f$ and $g$ have exponents $0$ or $1$ for each variable,
(ii) any member of the Gr\"obner basis of the ideal $(f,g)$ is not divided by any two-term factor, and
(iii) $f$ and $g$ become zero when $x=y=z=1$.
In the main text, the {\em cubic code} refers specifically to the code 1 in the table.
The numbering is consistent with a table in \cite{Haah2011Local}.}
\label{tb:10-models-wo-string}
\end{table}

\section{Cubic code}
\label{sec:cubic-code}


Written out explicitly, cubic code
is the translation-invariant negative sum of two types of interaction terms as in Figure~\ref{fig:CubicCode}.
%%%%%%%%%%%%%%%%%%%%%%%%%%%%%%
\begin{figure}[b]
\[
\drawgenerator{ZI}{ZZ}{IZ}{ZI}{IZ}{II}{ZI}{IZ} 
\quad
\drawgenerator{XI}{II}{IX}{XI}{IX}{XX}{XI}{IX}
\quad \quad
\drawgenerator{xz}{xyz}{x}{xy}{y}{1}{yz}{z}
\]
\caption{Stabilizer generators of the 3D cubic code. Here
$X\equiv \sigma^x$ and $Z\equiv \sigma^x$ represent single-qubit Pauli operators,
while $I$ is the identity operator.
Double-letter indices represent two-qubit Pauli operators,
for example,
$IZ\equiv I\otimes Z$, $ZZ\equiv Z\otimes Z$, $II\equiv I\otimes I$, etc.}
\label{fig:CubicCode}
\end{figure}
%%%%%%%%%%%%%%%%%%%%%%%%%%%%%%%
There are two qubits per site,
and the two-letter notation stands for tensor product of Pauli matrices.
For example, $XI = \sigma^x \otimes I$, $ZZ = \sigma^z \otimes \sigma^z$, etc.
The third cube specifies the coordinate system of the simple cubic lattice.
The generating map for the stabilizer module is
\[
\sigma_\text{cubic-code} =
\begin{pmatrix}
1 + x + y + z    & 0 \\
1 + xy + yz + zx & 0 \\
0 & 1 + \bar x \bar y + \bar y \bar z + \bar z \bar x \\
0 & 1 + \bar x + \bar y + \bar z \\
\end{pmatrix}
\]
where one has to interchange the first and second qubit.
The associated ideal is contained in a prime ideal of codimension 2 in $\FF_2[x^{\pm 1},y^{\pm 1},z^{\pm 1}]$:
\[
 I(\sigma) \subseteq ( 1+x+y+z,~ 1+xy+yz+zx ) = \pp_{xyz}.
\]
Since $\codim I(\sigma) \ge 2$, the characteristic dimension is 1.
Since $\coker \epsilon_\text{cubic-code} = R / \pp_{xyz} \oplus R / \overline{ \pp_{xyz}}$,
any nonzero element of $\pp_{xyz}$ is a fractal generator.
Since the conditions used in the search for the model
were only necessary conditions for the no-strings rule.
A rigorous treatments is as follows.
We prove a purely algebraic statement,
of which the no-strings rule is an interpretation.
A more elementary method can be found in \cite{Haah2011Local}.
%%%%%%%%%%%%%%%%%5
\begin{lem}
Let $S = \FF_2[x,y,z]$ be a polynomial ring, and $\pp = (1+x+y+z,~1+xy+yz+zx) \subseteq S$ an ideal.
If $m_1 e_1 + m_2 e_2 \in \pp$ for polynomials $e_1, e_2$
and monomials $m_1, m_2$ such that $\gcd( m_1, m_2 ) = 1$,
then only one of the following is true:
\begin{itemize}
\item $e_1 \in \pp$ and $e_2 \in \pp$.
\item $\max( \deg e_1, \deg e_2 ) \ge \max( \deg m_1, \deg m_2 )$.
\end{itemize}
\label{lem:alg-proof-no-string-cubic}
\end{lem}
\begin{proof}
Put $\FF_4 = \{ 0, 1, \omega, \omega^2 \}$,
i.e., $\omega$ is the primitive third root of unity over the binary field.
Consider a ring homomorphism $\phi : S \to U:=\FF_4[t]$ defined by
\[
\phi : x \mapsto 1+t, \quad y \mapsto 1+\omega t, \quad z \mapsto 1+\omega^2 t .
\]
It maps the two generators of $\pp$ to zero in $U$.
\begin{align*}
 \phi(1+x+y+z)    &= 4 + (1+\omega+\omega^2) t = 0 ,\\
 \phi(1+xy+yz+zx) &= 4 + 2(1+\omega+\omega^2)t + (1+\omega+\omega^2)t^2 = 0 .
\end{align*}
Moreover, $\ker \phi$ is precisely $\pp$.
(This can be verified by eliminating the variable $x$ 
and computing the Gr\"obner basis of $((y+1)+\omega(z+1),\omega^2+\omega+1)$ in an elimination monomial order.)
If $m_1 e_1 + m_2 e_2 \in \pp$, then $\phi(m_1)\phi(e_1) = \phi(m_2)\phi(e_2)$.
Since $m_1$ and $m_2$ are co-prime monomials and $\phi(x),\phi(y),\phi(z)$ are pairwise co-prime,
it follows that $\phi(m_1)$ and $\phi(m_2)$ are nonzero and co-prime.
Therefore, $\phi(e_1) = 0$ if and only if $\phi(e_2) = 0$ if and only if $e_1, e_2 \in \pp$,
which is the first case.
If $\phi(e_1) \neq 0$, then $\phi(m_1)$ must divide $\phi(e_2)$ and $\phi(m_2)$ must divide $\phi(e_1)$.
Since $\phi$ is degree-preserving, we have the second case.
\end{proof}
%%%%%%%%%%%%%%%%
\begin{theorem}
The cubic code obeys the no-strings rule with the constant $\alpha = 1$ under $\ell_\infty$-metric.
\label{thm:CubicCode-no-strings}
\end{theorem}
\begin{proof}
Since the cubic code is translationally invariant,
we may use the formalism of Chapter~\ref{chap:alg-theory}.
Since the cubic code is of CSS type, where $\coker \epsilon$ is a direct sum of isomorphic summands,
we only have to consider $\sigma^x$-type charges.
The set of all virtual charges is $R^1 = R = \FF_2[x^{\pm 1},y^{\pm 1},z^{\pm 1}]$
and the set of all trivial charges are given by a submodule (ideal) $\pp_{xyz} = (1+x+y+z,1+xy+yz+zx) \subseteq R$.
Note that $R$ is the localization of $S = \FF_2[x,y,z]$ by a single element $xyz$,
and $\pp_{xyz}$ is the localization of $\pp = (1+x+y+z,~1+xy+yz+zx)$ by the same single element $xyz$.

Let $e_1$ and $e_2$ be the charges contained in two boxes of a string segment.
Since they are overall trivial, we have $e_1 + x^i y^j z^k e_2 \in \pp$ where $i,j,k \in \ZZ$.
Equivalently, we may write $m_1 e_1 + m_2 e_2 \in \pp$ where $e_1$ and $e_2$ have nonnegative exponents,
and $m_1, m_2$ are monomials such that each variable ($x$, $y$, or $z$) appears
only in one of $m_1$ or $m_2$.
Since $\pp_{xyz} \cap S = \pp$, which can be verified using Gr\"obner basis techniques,
we are in the situation of Lemma~\ref{lem:alg-proof-no-string-cubic}.
The width of the string segment is the maximum of the degrees of $e_1$ and $e_2$,
and the $\ell_1$-distance between the boxes enclosing $e_1$ and $e_2$ is $\le \max( \deg m_1, \deg m_2 )$.
If we use $\ell_\infty$-distance, the constant $\alpha$ in the no-strings rule is $1$.
\end{proof}


Let us explicitly calculate the ground-state degeneracy
when the Hamiltonian is defined on $L \times L \times L$ 
cubic lattice with periodic boundary conditions.
By Corollary~\ref{cor:k-formulas},
\[
 k = \dim_{\FF_2} R / (\pp_{xyz} + \bb_L) \oplus R / (\overline{\pp_{xyz}} + \bb_L) = 2 \dim_{\FF_2} R / (\pp_{xyz} + \bb_L),
\]
where $\bb_L = ( x^L -1,~ y^L -1,~ z^L -1 )$.
So the calculation of ground-state degeneracy comes down to the calculation of
\[
 d = \dim_{\FF_2} T' / \pp
\]
where $T' = \FF_2[x,y,z]/(x^{n_1}-1, y^{n_2}-1, z^{n_3}-1 )$.

We may extend the scalar field to any extension field
without changing $d$.
Let $\FF$ be the algebraic closure of $\FF_2$ and let
\[
 T=\FF[x,y,z]/(x^{n_1}-1, y^{n_2}-1, z^{n_3}-1 )
\]
be an Artinian ring.
By Proposition~\ref{prop:Artin-ring},
it suffices to calculate for each maximal ideal $\mm$ of $T$
the vector space dimension 
\[
d_\mm = \dim_{\FF} (T / \pp)_\mm
\]
of the localized rings, and sum them up.

Suppose $n_1,n_2,n_3 > 1$.
By Nullstellensatz, any maximal ideal of $T$ is of form
$\mm=(x-x_0,y-y_0,z-z_0)$ where $x_0^{n_1}=y_0^{n_2}=z_0^{n_3}=1$.
(If $n_1 = n_2 = n_3 = 1$, then $T$ becomes a field, and there is no maximal ideal other than zero.)
Put $n_i = 2^{l_i}n_i'$ where $n_i'$ is not divisible by $2$.
Since the polynomial $x^{n_1} -1$ contains the factor $x-x_0$ with multiplicity $2^{l_1}$,
it follows that
\[
 T_\mm = \FF [x,y,z]_\mm / ( x^{2^{l_1}}+a',~ y^{2^{l_2}}+b',~ z^{2^{l_3}}+c')
\]
where $a' = x_0^{2^{l_1}}, b' = y_0^{2^{l_2}}, c' = z_0^{2^{l_3}}$.
Hence, $(T/\pp)_\mm \cong \FF[x,y,z] / I'$ where
\[
 I' = (x+y+z+1, xy+xz+yz+1 ,~ x^{2^{l_1}}+a',~ y^{2^{l_2}}+b',~ z^{2^{l_3}}+c').
\]
If $I' = \FF[x,y,z]$, then $d_\mm = 0$.


Without loss of generality, we assume that $l_1 \le l_2 \le l_3$.
By powering the first two generators of $I'$,
we see that $(x_0,y_0,z_0)$ must be a solution of them in order for $I'$ not to be a unit ideal.
Eliminating $z$ and shifting $x \to x+1$, $y \to y+1$,
our objective is to calculate the Gr\"obner basis for the proper ideal
\[
 I = (x^2+xy+y^2, x^{2^{l_1}}+a,~ y^{2^{l_2}}+b )
\]
where $a=a'+1$ and $b = b'+1$. So
\[
 d_\mm = \dim_{\FF} \FF[x,y]/I.
\]
One can easily deduce by induction that $y^{2^m} + x^{2^m-1}(m x +y) \in I$ for any integer $m \ge 0$.
And $b = \omega a^{2^{l_2 - l_1}}$ for a primitive third root of unity $\omega$.
So we arrive at
\[
 I = ( y^2 + yx + x^2,~ yx^{2^{l_2} -1 } + b ( 1 + l_2 \omega^2 ),~ x^{2^{l_1}} + a )
\]
We apply the Buchberger criterion.
If $a \neq 0$, i.e., $x_0 \neq 1$,
then $b \neq 0$ and $I = (x+(\omega^2 + l_2)y, x^{2^{l_1}} + a)$,
so $d_\mm = 2^{l_1}$

If $a=b=0$,
then $I = (y^2 + yx + x^2, yx^{2^{l_2}-1}, x^{2^{l_1}} )$.
The three generators form Gr\"obner basis if $l_2 = l_1$.
Thus, in this case, $d_\mm = 2^{l_1+1}-1$.
If $ l_2 > l_1$, then $d_\mm = 2^{l_1+1}$.

To summarize, except for the special point $(1,1,1)\in \FF^3$ of the affine space,
each point in the algebraic set
\[
 V = \left\{ (x,y,z) \in \FF^3 ~\middle|~
\begin{matrix}
x+y+z+1 = xy+xz+yz+1 = 0 \\  
x^{n_1'}-1 = y^{n_2'}-1 = z^{n_3'}-1 = 0
\end{matrix}
\right\}
\]
contribute $2^{l_1}$ to $d$. The contribution of $(1,1,1)$ is either $2^{l_1 +1}$ or $2^{l_1 +1}-1$.
The latter occurs if and only if $l_1$ and $l_2$,
the two smallest numbers of factors of $2$ in $n_1,n_2,n_3$,
are equal. Let $d_0 = \#V$ be the number of points in $V$.
The desired answer is
\[
 d = 2^{l_1} (d_0 -1) + 
\begin{cases} 
2^{l_1 +1} -1 & \text{if $l_1 = l_2$ } \\
2^{l_1 +1}    & \text{otherwise}
\end{cases}
\]
where $l_1 \le l_2 \le l_3$ are the number of factors of $2$ in $n_i$.


The algebraic set defined by $(x+y+z+1,~xy+xz+yz+1)$ is the union of two isomorphic lines
intersecting only at $x=y=z=1$, one of which is parametrized by $x \in \FF$ as
\[
 (1+x, 1+\omega x, 1+\omega^2 x) \in \FF^3,
\]
and another is parametrized as
\[
 (1+x,1+\omega^2 x, 1+\omega x) \in \FF^3.
\]
where $\omega$ is a primitive third root of unity.
Therefore, the purely geometric number $d_0 = 2 d_1 -1$ can be calculated by
\[
 d_1 = \deg_x \gcd \left( (1+x)^{n_1'}+1, (1+\omega x)^{n_2'}+1, (1+\omega^2 x)^{n_3'}+1 \right) .
\]
Using $(\alpha+\beta)^{2^p} = \alpha^{2^p} + \beta^{2^p}$ and $\omega^2 + \omega + 1 = 0$,
one can easily compute some special cases as summarized in the following corollary.
Some values of $k$ for small $L$ are presented in Table~\ref{tb:k-numerics} and Figure~\ref{fig:k-numerics}.

\begin{figure}[htbp]
\centering
\begin{minipage}{.45\textwidth}
\includegraphics[width=\textwidth]{k-numerics-graph-1}
\end{minipage}
\begin{minipage}{.45\textwidth}
\includegraphics[width=\textwidth]{k-numerics-graph-2}
\end{minipage}
\caption{Number of encoded qubits $k$ of the cubic code
defined on $L\times L \times L$ periodic lattice}
\label{fig:k-numerics}
\end{figure}


\begin{cor}
\label{cor:cubic-code-degeneracy-formula}
Let $2^k$ be the ground-state degeneracy of
the cubic code on the cubic lattice of size $L^3$ with periodic boundary conditions.
($k=k(L)$ is the number of encoded qubits.)
Then
\begin{align*}
\frac{k+2}{4} 
&= \deg_x \gcd \left( (1+x)^{L}+1,~ (1+\omega x)^{L}+1,~ (1+\omega^2 x)^{L}+1 \right)_{\FF_4} \\
&= \begin{cases}
     1    & \text{if $L = 2^p+1$}, \\
     L    & \text{if $L =2^p$}, \\
     L-2  & \text{if $L = 4^p -1$}, \\
     1    & \text{if $L = 2^{2p+1} -1$}.
   \end{cases}
\label{eq:cubic-code-formula-k}
\end{align*}
where $\omega^2 + \omega + 1 = 0$ and $p \ge 1$ is any integer.
If $L = 2^r L'$, then $k(L)+2 = 2^r( k(L') + 2 )$.
\end{cor}

\begin{table}[tp]
\centering
\begin{tabular}{r|rll}
\hline
\multicolumn{4}{c}{$(k+2)/4 = 1 + 12 \sum_n q_n(L)$} \\
\multicolumn{4}{c}{$q_n(L)$ is nonzero only if $n | L$.} \\
\hline
$q_n(L)$&  & $n$ & \\
\hline
1 & 15 & $= 2^4 -1$ & $=3 \cdot 5$  \\
5 & 63 & $= 2^6 -1$ & $=3^2 \cdot 7$  \\
20 & 255 & $= 2^8 -1$ & $=3 \cdot 5 \cdot 17$  \\
80 & 1023 & $=2^{10}-1$ & $=3 \cdot 11 \cdot 31$  \\
322 & 4095 & $= 2^{12} -1$ & $=3^2 \cdot 5 \cdot 7 \cdot 13$ \\
\hline
5 & 341 & $=(2^{10}-1)/(2^2 -1)$ & $= 11 \cdot 31$  \\
6 & 1365 & $=(2^{12}-1)/(2^2 -1)$ & $=3 \cdot 5 \cdot 7 \cdot 13$  \\
49 & 5461 & $=(2^{14}-1)/(2^2 -1)$ & $=43 \cdot 127$  \\
\hline
4 & 455 & $=(2^{12}-1)/(2^3 +1)$ & $= 5 \cdot 7 \cdot 13$  \\
3 & 585 & $=(2^{12}-1)/(2^3 -1)$ & $= 3^2 \cdot 5 \cdot 13$  \\
9 & 9709& $=(2^{18}-1)/(3(2^3+1))$& $= 7\cdot 19 \cdot 73$  \\
5 & 11275&$=11(2^{10}+1)$ & $=5^2 \cdot 11 \cdot 41$ \\
\hline
\end{tabular}
\caption{Numerical values of $k$ for \emph{odd} linear size $L$
computed from Corollary~\ref{cor:cubic-code-degeneracy-formula}.
The list is complete if $2 \le L \le 20000$.
For example, if $L=945=15 \cdot 63$, then $\frac{k+2}{4} = 1 + 12\cdot(1+5) = 73$.
}
\label{tb:k-numerics}
\end{table}


\section{Real-space renormalization of the cubic code}
\label{sec:real-space-RG-CubicCode}

Renormalization group refers to a machinery 
to extract essential properties of the system at long distances.
For lattice spin models, the so-called ``block spin'' method amounts to
considering a sequence of coarse-grained lattices
and finding effective Hamiltonians pertaining to the coarse-grained lattices~\cite{Wilson1975RG,Kadanoff1977RG}.
In the sequence of coarse-graining, short-ranged correlations will disappear,
and long-ranged essential correlations will remain.
Thus, if a generic Hamiltonian retains its form under the renormalization group flow,
we expect that the Hamiltonian should represent a proper phase of matter,
and can wonder about universal aspects of the phase.

For topologically ordered systems,
it is interesting to look at entanglement structure instead of correlation functions,
since, as they are gapped, the correlation functions of local operators would decay 
exponentially with the distance between the regions the operators act on~\cite{HastingsKoma2006CorDecay},
and hence, if they are renormalization group fixed points, they will have zero correlation length.
As we wish to ignore local deformations of the system,
it is legitimate to apply a finite depth quantum circuit (local unitary) to the system
so as to remove some local
entanglement~\cite{VerstraeteCiracLatorreEtAl2005Renormalization,
Vidal2007ER,AguadoVidal2007Entanglement,Vidal2008class,ChenGuWen2010transformation}.
More precisely, one applies a finite composition of unitary operators on a ground state,
each of which can be written as a product of local unitary operators of disjoint supports,
and then try to identify spins in a product state.

Code Hamiltonians admit an even simpler renormalization scheme.
Under the local unitaries the Hamiltonian is conjugated,
and spins in product states is readily identified by single-spin operators
such as $\sigma^z$ acting on unentangled spins.
If we restrict ourselves to translationally invariant case,
we may also assume that the local unitaries obey the translation invariance.
For example, it is well-known that the 2D toric code model
is a renormalization group fixed point under this scheme%
~\cite{AguadoVidal2007Entanglement}.

In Chapter~\ref{chap:alg-theory}, we have developed enough tools for the simplified renormalization group flow computation.
In fact, we have done a computation in Example~\ref{eg:wen-plaquette}.
To warm up, let us compute the renormalization of the 2D toric code.
The generating matrix is
\[
 \sigma =
\begin{pmatrix}
 1+x & 0 \\
 1+y & 0 \\
\hline
\hline
  0  & 1+\frac{1}{y} \\
  0  & 1+\frac{1}{x} \\
\end{pmatrix}
\]
If we coarse-grain the lattice by blocking two sites in $x$-direction,
then every module is now viewed as a module over $\FF_2[x^{\pm 2}, y^{\pm 1}]$,
and each entry of $\sigma$ is replaced by a $2 \times 2$ matrix,
since $\sigma$ is a map between free modules $G$ and $P$.
Concretely,
\[
 \sigma' = 
\begin{pmatrix}
 1 & 1 & 0 & 0 \\
 x & 1 & 0 & 0 \\
 1+y & 0 & 0 & 0 \\
 0 & 1+y & 0 & 0 \\
\hline
\hline
 0 & 0 & 1+\frac{1}{y} & 0 \\
 0 & 0 & 0 & 1+\frac{1}{y} \\
 0 & 0 & 1 & \frac{1}{x} \\
 0 & 0 & 1 & 1 
\end{pmatrix}
\]
where the double line distinguishes $\sigma^x$-part and $\sigma^z$-part.
Here, the variable $x$ really means the translation along $x$-direction by \emph{two} units of the original lattice.
Applying local unitaries, we see
\[
 \sigma' \to
\left(
\begin{array}{cccc}
 1 & 1 & 0 & 0 \\
 0 & 1+x & 0 & 0 \\
 0 & 1+y & 0 & 0 \\
 0 & 1+y & 0 & 0 \\
\hline
\hline
 0 & 0 & 0 & 0 \\
 0 & 0 & 0 & 1+\frac{1}{y} \\
 0 & 0 & 1 & \frac{1}{x} \\
 0 & 0 & 1 & 1
\end{array}
\right)
\to
\left(
\begin{array}{cccc}
 1 & 1 & 0 & 0 \\
 0 & 1+x & 0 & 0 \\
 0 & 1+y & 0 & 0 \\
 0 & 0 & 0 & 0 \\
\hline
\hline
 0 & 0 & 0 & 0 \\
 0 & 0 & 0 & 1+\frac{1}{y} \\
 0 & 0 & 0 & 1+\frac{1}{x} \\
 0 & 0 & 1 & 1
\end{array}
\right)
\to
\left(
\begin{array}{c|c|c|c}
 1 & 0 & 0 & 0 \\
\hline
 0 & 1+x & 0 & 0 \\
 0 & 1+y & 0 & 0 \\
\hline
 0 & 0 & 0 & 0 \\
\hline
\hline
 0 & 0 & 0 & 0 \\
\hline
 0 & 0 & 0 & 1+\frac{1}{y} \\
 0 & 0 & 0 & 1+\frac{1}{x} \\
\hline
 0 & 0 & 1 & 0
\end{array}
\right).
\]
In the last matrix, it is evident that the first qubit (all the first qubits on every site)
and the fourth qubit are disentangled.
Ignoring those, we see that the generating matrix and also the Hamiltonian retain the original form.
A drawback of this computation is that it is hard to understand why this should happen.
Fortunately, there is a better understanding for two-dimensional quantum double models~\cite{Kitaev2003Fault-tolerant},
which includes the 2D toric code model,
using so-called $G$-injective PEPS by Schuch, Cirac, and P\'erez-Garc\'ia~\cite{SchuchCiracPerez-Garcia2010G-injective}.
Their conclusion is that
the quantum double model is constructed with the regular representation of a finite symmetry group $G$,
and it is a renormalization group fixed point
because of the plethysm $\mathbb{C}G \otimes \mathbb{C}G \cong \mathbb{C}G \otimes (\mathbb{C}1)^{\oplus |G|}$ 
of the regular representation.

Now we return to the cubic code. We perform a similar computation as above,
and find that under blocking of $2 \times 2 \times 2$ sites (16 spins in total)
the cubic code model $A$ decomposes into two non-interacting Hamiltonians $A$ and $B$
living in the coarse-grained lattice,
one of which is the same as the original $A$,
but the other $B$ looks different.
Detailed calculation will be given below.
The generating matrices are as follows.
\[
 \sigma_A = 
\begin{pmatrix}
 1+x+y+z & 0 \\
 1+x y+y z+z x & 0 \\
\hline
\hline
 0 & 1+\frac{1}{xy}+\frac{1}{yz}+\frac{1}{zx} \\
 0 & 1+\frac{1}{x}+\frac{1}{y}+\frac{1}{z} \\
\end{pmatrix}, \quad
 \sigma_B = 
\left(
 \begin{array}{cccc}
 x+z & 1+x &   &   \\
 1+x & 1+z &   &   \\
 x+y & 1+y &   &   \\
 1+y & 1+x &   &   \\
\hline
\hline
   &   & 1+\frac{1}{y} & \frac{1}{x}+\frac{1}{y} \\
   &   & 1+\frac{1}{x} & 1+\frac{1}{y} \\
   &   & 1+\frac{1}{x} & \frac{1}{x}+\frac{1}{z} \\
   &   & 1+\frac{1}{z} & 1+\frac{1}{x}
\end{array}
\right).
\]
We can repeat the renormalization group flow computation for the model $B$ only.
We find that after blocking of $2 \times 2 \times 2$ sites,
$B$ is renormalized to the identical two copies of $B$ itself.
(Calculation will be given below.)
\[
 A \xrightarrow{2\times 2 \times 2} A \oplus B, \quad \text{ and } \quad
B \xrightarrow{2 \times 2 \times 2} B \oplus B.
\]
The renormalization group yields an explicit method to produce the ground states of the cubic code,
Start with a state one wish to encode, put auxiliary qubits in the trivial state,
apply the inverse local unitaries, and iterate the overall process on the \emph{refined} lattice with more auxiliary qubits.
It is slightly different from the MERA (Multiscale Entanglement Renormalization Ansatz)
prescription of 2D toric code state developed 
in~\cite{AguadoVidal2007Entanglement}.
Rather a so-called branching MERA~\cite{EvenblyVidal2012class} is more appropriate.

We do not have a deeper understanding why this should happen.
However, there are some consistency checks
from the degeneracy formula and the annihilator of the module of topological charges.
Corollary~\ref{cor:cubic-code-degeneracy-formula} says that
the number of encoded qubits $k$ is
\begin{align*}
k(L) = 4L -2 
&= \left( 4 \frac{L}{2} - 2\right) & &+ \left( 4 \frac{L}{2} \right) \\
&= \left( 4 \frac{L}{4} - 2\right) + \left( 4 \frac{L}{4} \right) & &+ \left( 4 \frac{L}{4} \right) + \left( 4 \frac{L}{4} \right)
\end{align*}
when the linear system size $L$ is a power of 2.
The expression is decomposed to display contributions from the model $A$ and the model $B$ explicitly.
It suggests that the model $A$, the original cubic code, 
cannot give rise to an identical pair of models at coarse-grained lattice,
whereas $B$ can.
It is not too clear whether the two models $A$ and $B$ are really non-isomorphic.

On the other hand, the annihilator of the topological charge module 
tells us there is something special about $2 \times 2 \times 2$ blocking.
The topological charge module is the torsion part of the virtual excitation module factored by trivial charge module,
i.e., the torsion submodule of $\coker \epsilon$.
As for the cubic code, $\coker \epsilon$ is a torsion module, which is a direct sum of two isomorphic summands.
\begin{align*}
& \coker \epsilon = R/\pp_{xyz} \oplus R / \overline{ \pp_{xyz} } , \\
& R = \FF_2[x^{\pm 1},y^{\pm 1},z^{\pm 1}],\\
& \pp_{xyz} = (1+x+y+z,1+xy+yz+zx) \subset R.\\
\end{align*}
where the bar denotes antipode map.
We focus on the first summand $R/\pp_{xyz}$.
If we replace $\FF_2$ with its algebraic closure $\FF$ for convenience,
$R/\pp_{xyz} = (\FF[x,y,z]/\pp)_{xyz}$ is the coordinate ring of an affine variety defined by $\pp$
localized on the complement of the union of three planes $xyz=0$.
The variety is a union of two isomorphic lines, 
as we have seen in the degeneracy calculation of Section~\ref{sec:cubic-code}.
One line is parameterized as
\[
x = 1+ t, \quad y = 1+ \omega t, \quad z = 1 + \omega^2 t
\]
and the other is
\[
x = 1+t, \quad y = 1+ \omega^2 t, \quad z = 1+ \omega t
\]
where $\omega$ is the primitive third root of unity.
The ideal $\pp_{xyz}$ is the annihilator of a module $M = R/\pp_{xyz}$.
Under coarse-graining $M$ is promoted to a module over the coarse-grained
translation group algebra, which is $R ' = \FF_2[x^{\pm 2}, y^{\pm 2}, z^{\pm 2}]$
in our $2 \times 2 \times 2$ blocking.
Then, the annihilator of $M$ in the new ring $R'$ is just $R' \cap I$.
It is easy to verify that $R' \cap \pp_{xyz} = (1+x^2 + y^2 + z^2, 1+ x^2 y^2 + y^2 z^2 + z^2 x^2 ) = \pp'$
using Gr\"obner basis.
As rings, $R'/\pp'$ and $R/\pp_{xyz}$ are isomorphic.
Thus, it is consistent that $A$ renormalizes to something similar to itself.
Further, we see that the charge annihilator of $B$ must be the same as that of $A$.

This observation tells us that $2^3$ blocking is special.
If we had blocked $3^3$ sites, we would not see the self-reproducing behavior,
since the charge annihilator would be different:
\begin{align*}
 \FF_2[x^{\pm 3}, y^{\pm 3}, z^{\pm 3}] \cap \pp_{xyz} = 
& \left( 1+{y'}+{y'}^2+{y'}^3+{z'}+{y'} {z'}+{y'}^2 {z'}+{z'}^2+{y'} {z'}^2+{z'}^3, \right. \\
& \left. 1+{x'}+{x'}^2+{y'}+{x'} {y'}+{y'}^2+{z'}+{x'} {z'}+{y'} {z'}+{z'}^2 \right)
\end{align*}
where $x'=x^3,y'=y^3,z'=z^3$.

\subsection*{Calculation}

The generating matrix $\sigma_A$
for the stabilizer module of the cubic code
transforms under the coarse-graining by blocking two sites along $x$-direction as
\begin{align*}
 \sigma_A \xrightarrow{\text{coarse-grain }x^2 \to x} \quad
& \sigma_1 = \begin{pmatrix} \sigma_{1X} & 0 \\ 0 & \sigma_{1Z} \end{pmatrix} \\
& \sigma_{1X} =
\begin{pmatrix}
 1+y+z   & 1     \\
 x       & 1+y+z \\
 1+y z   & y+z   \\
 x y+x z & 1+y z \\
\end{pmatrix}, \quad
\sigma_{1Z} =
\begin{pmatrix}
 1+\frac{1}{y z} & \frac{1}{x y}+\frac{1}{x z} \\
 \frac{1}{y}+\frac{1}{z} & 1+\frac{1}{y z} \\
 1+\frac{1}{y}+\frac{1}{z} & \frac{1}{x} \\
 1 & 1+\frac{1}{y}+\frac{1}{z} \\
\end{pmatrix}.
\end{align*}
We apply elementary symplectic transformations.
Recall $\dagger$ is the transpose followed by entry-wise antipode 
map $x \mapsto x^{-1}$, $y \mapsto y^{-1}$, $z \mapsto z^{-1}$.
\[
\sigma_2 = \begin{pmatrix} \sigma_{2X} & 0 \\ 0 & \sigma_{2Z} \end{pmatrix}
 = 
\begin{pmatrix} r_1 & 0 \\ 0 & r_1^\dagger \end{pmatrix} 
\begin{pmatrix} \sigma_{1X} & 0 \\ 0 & \sigma_{1Z} \end{pmatrix}
\begin{pmatrix} c_1 & 0 \\ 0 & c_1^\dagger \end{pmatrix}
 \]
where
\begin{align*}
\sigma_{2X} &=
\left(\begin{array}{c|c}
 0 & 1 \\
\hline
 1+x+y^2+z^2 & 0  \\
 1+y+y^2+z+y z+z^2 & 0  \\
\hline
 0 & 0  \\
\end{array}\right),
%%%%
& &
%%%%
\sigma_{2Z} =
\left(
\begin{array}{c|c}
 0 & 0 \\
\hline
0 & 1+\frac{1}{y^2}+\frac{1}{y}+\frac{1}{z^2}+\frac{1}{z}+\frac{1}{y z} \\
0 & 1+\frac{1}{x}+\frac{1}{y^2}+\frac{1}{z^2} \\
\hline
1 & 0 \\
\end{array}
\right) ,
%%%%%%
\\
%%%%%%
 r_1 &= 
\begin{pmatrix}
 1     & 0   & 0     & 0 \\
 1+y+z & 1   & 0     & 0 \\
 y+z   & 0   & 1     & 0 \\  
 1+y z & y+z & 1+y+z & 1 \\
\end{pmatrix},
%%%%%%
& &
%%%%%
c_1 =
\begin{pmatrix}
 1     & 0 \\
 1+y+z & 1 \\
\end{pmatrix}.
\end{align*}
The first and fourth qubit may be factored out from $\sigma_2$.
A subsequent coarse-graining by blocking two sites in each $y$- and $z$-direction gives
\begin{align*}
\sigma_2 \xrightarrow{y^2 \to y,~~z^2 \to z} \quad
& \sigma_3 = \begin{pmatrix} \sigma_{3X} & 0 \\ 0 & \sigma_{3Z} \end{pmatrix} \\
\end{align*}
\begin{align*}
& \sigma_{3X} =
\begin{pmatrix}
 1+x+y+z & 0 & 0 & 0 \\
 0 & 1+x+y+z & 0 & 0 \\
 0 & 0 & 1+x+y+z & 0 \\
 0 & 0 & 0 & 1+x+y+z \\
 1+y+z & 1 & 1 & 1 \\
 z & 1+y+z & z & 1 \\
 y & y & 1+y+z & 1 \\
 y z & y & z & 1+y+z
\end{pmatrix},
\end{align*}
\begin{align*}
 & \sigma_{3Z} =
\begin{pmatrix}
 1+\frac{1}{y}+\frac{1}{z} & \frac{1}{z} & \frac{1}{y} & \frac{1}{y z} \\
 1 & 1+\frac{1}{y}+\frac{1}{z} & \frac{1}{y} & \frac{1}{y} \\
 1 & \frac{1}{z} & 1+\frac{1}{y}+\frac{1}{z} & \frac{1}{z} \\
 1 & 1 & 1 & 1+\frac{1}{y}+\frac{1}{z} \\
 1+\frac{1}{x}+\frac{1}{y}+\frac{1}{z} & 0 & 0 & 0 \\
 0 & 1+\frac{1}{x}+\frac{1}{y}+\frac{1}{z} & 0 & 0 \\
 0 & 0 & 1+\frac{1}{x}+\frac{1}{y}+\frac{1}{z} & 0 \\
 0 & 0 & 0 & 1+\frac{1}{x}+\frac{1}{y}+\frac{1}{z} \\
\end{pmatrix}
\end{align*}
We keep applying elementary symplectic transformations.
\[
 \sigma_4 = \begin{pmatrix} \sigma_{4X} & 0 \\ 0 & \sigma_{4Z} \end{pmatrix}
=
\begin{pmatrix} r_{2X} & 0 \\ 0 & r_{2Z} \end{pmatrix}
\begin{pmatrix} \sigma_{3X} & 0 \\ 0 & \sigma_{3Z} \end{pmatrix}
\begin{pmatrix} c_{2X} & 0 \\ 0 & c_{2Z} \end{pmatrix}
\]
where
\begin{align*}
\sigma_{4X} = 
\left(\begin{array}{c|c|cc}
 0 & 1 & 0 & 0  \\
 0 & 0 & 0 & 0  \\
\hline
 0 & 0 & 1+x+y+z & 0  \\
 0 & 0 & 0 & 1+x+y+z  \\
 1+x+y+z & 0 & 0 & 0  \\
 x+x y+z+x z & 0 & 1+y & y+z \\
 y+x y & 0 & 1+z & 1+y       \\
 x y+y z & 0 & y+z & 1+z     \\
\end{array}\right),
\end{align*}
\begin{align*}
\sigma_{4Z} =
\left(\begin{array}{c|ccc}
 0 & 0 & 0 & 0 \\
 1 & 0 & 0 & 0 \\
\hline
 0 & 1+\frac{1}{y} & 1+\frac{1}{z} & \frac{1}{y}+\frac{1}{z} \\
 0 & \frac{1}{y}+\frac{1}{z} & 1+\frac{1}{y} & 1+\frac{1}{z} \\
 0 & \frac{1}{x}+\frac{1}{x y}+\frac{1}{z}+\frac{1}{x z} & \frac{1}{y}+\frac{1}{x y} & \frac{1}{x y}+\frac{1}{y z} \\
 0 & 1+\frac{1}{x}+\frac{1}{y}+\frac{1}{z} & 0 & 0 \\
 0 & 0 & 1+\frac{1}{x}+\frac{1}{y}+\frac{1}{z} & 0 \\
 0 & 0 & 0 & 1+\frac{1}{x}+\frac{1}{y}+\frac{1}{z}
\end{array}\right),
\end{align*}
\begin{align*}
r_{2X} = 
\begin{pmatrix}
 1 & 0 & 0 & 0 & 1 & 0 & 0 & 0   \\
 1+y+z & 1 & 1 & 1 & 1+x+y+z & 0 & 0 & 0   \\
 0 & 0 & 1 & 0 & 0 & 0 & 0 & 0   \\
 0 & 0 & 0 & 1 & 0 & 0 & 0 & 0   \\
 1 & 0 & 0 & 0 & 0 & 0 & 0 & 0   \\
 1+y+z & 0 & 0 & 0 & 1+y+z & 1 & 0 & 0   \\
 y & 0 & 0 & 0 & y & 0 & 1 & 0   \\
 y & 0 & 0 & 0 & y & 0 & 0 & 1   \\
\end{pmatrix},
\end{align*}
\begin{align*}
r_{2Z} =
\begin{pmatrix}
 0 & 1+\frac{1}{x}+\frac{1}{y}+\frac{1}{z} & 0 & 0 & 1 & 1+\frac{1}{y}+\frac{1}{z} & \frac{1}{y} & \frac{1}{y} \\
 0 & 1 & 0 & 0 & 0 & 0 & 0 & 0 \\
 0 & 1 & 1 & 0 & 0 & 0 & 0 & 0 \\
 0 & 1 & 0 & 1 & 0 & 0 & 0 & 0 \\
 1 & \frac{1}{x} & 0 & 0 & 1 & 0 & 0 & 0 \\
 0 & 0 & 0 & 0 & 0 & 1 & 0 & 0 \\
 0 & 0 & 0 & 0 & 0 & 0 & 1 & 0 \\
 0 & 0 & 0 & 0 & 0 & 0 & 0 & 1 \\
\end{pmatrix}.
\end{align*}
\begin{align*}
c_{2X} =
\begin{pmatrix}
 1 & 0 & 0 & 0 \\
 x & 1 & 1 & 1 \\
 0 & 0 & 1 & 0 \\
 0 & 0 & 0 & 1 \\
\end{pmatrix}, \quad
c_{2Z} =
\begin{pmatrix}
 1 & 1+\frac{1}{y}+\frac{1}{z} & \frac{1}{y} & \frac{1}{y} \\
 0 & 1 & 0 & 0 \\
 0 & 0 & 1 & 0 \\
 0 & 0 & 0 & 1
\end{pmatrix},
\end{align*}
Factoring out the first and second qubit from $\sigma_4$
and applying elementary symplectic transformations,
we arrive at
\[
 \sigma_5 = 
\begin{pmatrix} \sigma_{5X} & 0 \\ 0 & \sigma_{5Z} \end{pmatrix}
=
\begin{pmatrix} r_{3X} & 0 \\ 0 & r_{3Z} \end{pmatrix}
\left( \begin{pmatrix} \sigma_{4X} & 0 \\ 0 & \sigma_{4Z} \end{pmatrix} \rvert_{\{1,2\}^c} \right)
\begin{pmatrix} c_{3X} & 0 \\ 0 & c_{3Z} \end{pmatrix}
\]
where
\begin{align*}
\sigma_{5X} =
\begin{pmatrix}
 1+x+y+z & 0 & 0  \\
 1+x y+x z+y z & 0 & 0  \\
 0 & x+z & 1+x  \\
 0 & 1+x & 1+z  \\
 0 & x+y & 1+y  \\
 0 & 1+y & 1+x  \\
\end{pmatrix},
%%%%%
\quad
%%%%%
\sigma_{5Z} =
\begin{pmatrix}
 0 & 0 & 1+\frac{1}{x y}+\frac{1}{x z}+\frac{1}{y z} \\
 0 & 0 & 1+\frac{1}{x}+\frac{1}{y}+\frac{1}{z} \\
 1+\frac{1}{y} & \frac{1}{x}+\frac{1}{y} & 0 \\
 1+\frac{1}{x} & 1+\frac{1}{y} & 0 \\
 1+\frac{1}{x} & \frac{1}{x}+\frac{1}{z} & 0 \\
 1+\frac{1}{z} & 1+\frac{1}{x} & 0
\end{pmatrix},
\end{align*}
\begin{align*}
r_3 =
\begin{pmatrix}
 0 & 0 & 1 & 0 & 0 & 0 \\
 0 & 0 & 1 & 1 & 1 & 1 \\
 1 & 1 & x & 1 & 0 & 0 \\
 1 & 0 & 1 & 1 & 1 & 0 \\
 1 & 0 & 0 & 0 & 1 & 0 \\
 0 & 1 & 1+x & 1 & 0 & 0  \\
\end{pmatrix},
\quad
r_{3Z} =
\begin{pmatrix}
1 & \frac{1}{x} & 1 & 1 & 1 & 1 \\
0 & 0 & 0 & 0 & 0 & 1 \\
1 & 0 & 0 & 0 & 1 & 1 \\
0 & 1 & 0 & 1 & 0 & 1 \\
0 & 1 & 0 & 1 & 1 & 0 \\
1 & 1 & 0 & 0 & 1 & 1
\end{pmatrix},
\end{align*}
\begin{align*}
c_3 =
\begin{pmatrix}
 1 & 0 & 0 \\
 1 & 1 & 0 \\
 x & 0 & 1 \\
\end{pmatrix},
\quad
c_{3Z} =
\begin{pmatrix}
1 & 0 & 1 \\
0 & 1 & 1 \\
0 & 0 & 1
\end{pmatrix}.
\end{align*}
It is clear that $\sigma_A$, $\sigma_1$, $\sigma_2$, $\sigma_3$, $\sigma_4$, and $\sigma_5$ are all equivalent
because we applied only elementary symplectic transformations.
$\sigma_5$ shows a decomposition of $\sigma$ into two non-interacting two models,
one of which is $\sigma_A$ at $2 \times 2 \times 2$ coarse-grained lattice and another is $\sigma_B$.
%%%%%
We perform a similar renormalization for $\sigma_B$.
The coarse-graining by blocking two sites in $x$-direction gives
\begin{align*}
\sigma_B \xrightarrow{x^2 \to x} \quad
&\sigma_{B1} = \begin{pmatrix} \sigma_{B1X} & 0 \\ 0 & \sigma_{B1Z} \end{pmatrix} \\
&
\sigma_{B1X} =
\begin{pmatrix}
 z & 1 & 1 & 1   \\
 x & z & x & 1   \\
 1 & 1 & 1+z & 0 \\
 x & 1 & 0 & 1+z \\
 y & 1 & 1+y & 0 \\
 x & y & 0 & 1+y \\
 1+y & 0 & 1 & 1 \\
 0 & 1+y & x & 1 \\
\end{pmatrix},
\quad
\sigma_{B1Z} =
\begin{pmatrix}
 1+\frac{1}{y} & 0 & \frac{1}{y} & \frac{1}{x} \\
 0 & 1+\frac{1}{y} & 1 & \frac{1}{y} \\
 1 & \frac{1}{x} & 1+\frac{1}{y} & 0 \\
 1 & 1 & 0 & 1+\frac{1}{y} \\
 1 & \frac{1}{x} & \frac{1}{z} & \frac{1}{x} \\
 1 & 1 & 1 & \frac{1}{z} \\
 1+\frac{1}{z} & 0 & 1 & \frac{1}{x} \\
 0 & 1+\frac{1}{z} & 1 & 1
\end{pmatrix}.
\end{align*}
Apply elementary symplectic transformations to factor out trivial qubits.
\[
\sigma_{B2} =
 \begin{pmatrix} \sigma_{B2X} & 0 \\ 0 & \sigma_{B2Z} \end{pmatrix}
=
\begin{pmatrix} r_{4X} & 0 \\ 0 & r_{4Z} \end{pmatrix}
\begin{pmatrix} \sigma_{B1X} & 0 \\ 0 & \sigma_{B1Z} \end{pmatrix}
\begin{pmatrix} c_{4X} & 0 \\ 0 & c_{4Z} \end{pmatrix}
\]
where
{\small
\begin{align*}
&\sigma_{B2X} =                    & & \sigma_{B2Z} = \\
&\left(
\begin{array}{c|cc|c}
 0 & 0 & 0 & 1 \\
 0 & 0 & 0 & 0 \\
 1 & 0 & 0 & 0 \\
 0 & 0 & 0 & 0 \\
\hline
 0 & 1+y & y+z & 0       \\
 0 & 1+x+z+y z & y+z & 0 \\
 0 & 0 & x+y+z+y z & 0   \\
 0 & y+z & 1+x+y+y z & 0 \\
\end{array}
\right),
& &
%\sigma_{B2Z} =
\left(\begin{array}{c|c|c|c}
 0 & 0 & 0 & 0 \\
 0 & 0 & 1 & 0 \\
 0 & 0 & 0 & 0 \\
 1 & 0 & 0 & 0 \\
\hline
 0 & 1+\frac{1}{x}+\frac{1}{z}+\frac{1}{y z} & 0 & \frac{1}{y}+\frac{1}{z} \\
 0 & 1+\frac{1}{y} & 0 & \frac{1}{y}+\frac{1}{z} \\
 0 & \frac{1}{y}+\frac{1}{z} & 0 & 1+\frac{1}{x}+\frac{1}{y}+\frac{1}{y z} \\
 0 & 0 & 0 & \frac{1}{x}+\frac{1}{y}+\frac{1}{z}+\frac{1}{y z}
\end{array}\right),
\end{align*}
\begin{align*}
& r_{4X} = &  & r_{4Z} = \\
& \begin{pmatrix}
 1 & 0 & z & 0 & 0 & 0 & 0 & 0    \\
 y & 1 & 1+y & 0 & z & 1 & 1 & 1  \\
 0 & 0 & 1 & 0 & 0 & 0 & 0 & 0    \\
 1+y & 0 & 1 & 1 & 1 & 1 & 1+z & 0\\
 0 & 0 & y & 0 & 1 & 0 & 0 & 0         \\
 1+y & 0 & x+z+y z & 0 & 0 & 1 & 0 & 0 \\
 0 & 0 & 1+y & 0 & 0 & 0 & 1 & 1       \\
 1 & 0 & z & 0 & 0 & 0 & 0 & 1         \\
\end{pmatrix},
& &
% r_{4Z} =
\begin{pmatrix}
 1 & 1 & 0 & 1+\frac{1}{z} & 0 & 1+\frac{1}{y} & 1 & 1 \\
 0 & 1 & 0 & 0 & 0 & 0 & 0 & 0 \\
 \frac{1}{z} & \frac{1}{x} & 1 & \frac{1}{x} & \frac{1}{y} & \frac{1}{x} & 1+\frac{1}{y} & 0 \\
 0 & 0 & 0 & 1 & 0 & 0 & 0 & 0 \\
 0 & \frac{1}{z} & 0 & 1 & 1 & 0 & 0 & 0 \\
 0 & 1 & 0 & 1 & 0 & 1 & 0 & 0 \\
 0 & 1 & 0 & 1+\frac{1}{z} & 0 & 0 & 1 & 0 \\
 0 & 0 & 0 & 1+\frac{1}{z} & 0 & 0 & 1 & 1 \\
\end{pmatrix}.
\end{align*}
\[
c_{4X} =
\begin{pmatrix}
 1 & 1 & 0 & 0   \\
 0 & 1 & 1+z & 0 \\
 0 & 0 & 1 & 0   \\
 0 & 1+z & z & 1 \\
\end{pmatrix},
\quad
c_{4Z} =
\begin{pmatrix}
1 & 1 & 0 & \frac{1}{y} \\
0 & 1 & 0 & 1 \\
0 & 1+\frac{1}{y} & 1 & 1 \\
0 & 0 & 0 & 1
\end{pmatrix},
\]
}%small font
Factoring out trivial qubits from $\sigma_{B2}$ and coarse-graining by blocking two sites along $z$-direction,
we have
\begin{align*}
& \sigma_{B2}|_{\{1,2,3,4\}^c} \xrightarrow{z^2 \to z} \quad 
\sigma_{B3} = \begin{pmatrix} \sigma_{B3X} & 0 \\ 0 & \sigma_{B3Z} \end{pmatrix}, \text{ where }
\end{align*}
\begin{align*}
& \sigma_{B3X} = & & \sigma_{B3Z} =\\
& \begin{pmatrix}
 1+y & 0 & y & 1    \\
 0 & 1+y & z & y    \\
 1+x & 1+y & y & 1  \\
 z+y z & 1+x & z & y\\
 0 & 0 & x+y & 1+y  \\
 0 & 0 & z+y z & x+y\\
 y & 1 & 1+x+y & y  \\
 z & y & y z & 1+x+y\\
\end{pmatrix},
& &
% \sigma_{B3Z} =
\begin{pmatrix}
 1+\frac{1}{x} & \frac{1}{z}+\frac{1}{y z} & \frac{1}{y} & \frac{1}{z} \\
 1+\frac{1}{y} & 1+\frac{1}{x} & 1 & \frac{1}{y} \\
 1+\frac{1}{y} & 0 & \frac{1}{y} & \frac{1}{z} \\
 0 & 1+\frac{1}{y} & 1 & \frac{1}{y} \\
 \frac{1}{y} & \frac{1}{z} & 1+\frac{1}{x}+\frac{1}{y} & \frac{1}{y z} \\
 1 & \frac{1}{y} & \frac{1}{y} & 1+\frac{1}{x}+\frac{1}{y} \\
 0 & 0 & \frac{1}{x}+\frac{1}{y} & \frac{1}{z}+\frac{1}{y z} \\
 0 & 0 & 1+\frac{1}{y} & \frac{1}{x}+\frac{1}{y}
\end{pmatrix}.
\end{align*}
Again apply elementary symplectic transformations.
\[
\sigma_{B4} =
 \begin{pmatrix} \sigma_{B4X} & 0 \\ 0 & \sigma_{B4Z} \end{pmatrix}
=
\begin{pmatrix} r_{5X} & 0 \\ 0 & r_{5Z} \end{pmatrix}
\begin{pmatrix} \sigma_{B3X} & 0 \\ 0 & \sigma_{B3Z} \end{pmatrix}
\begin{pmatrix} c_{5X} & 0 \\ 0 & c_{5Z} \end{pmatrix}
\]
where
\[
 \sigma_{B4X} =
\left(\begin{array}{c|c|c|c}
 0 & 0 & 1 & 0 \\
\hline
 1+x+y+x y^2+z+y z & 0 & 0 & x y+y^2+x y^2+y z \\
\hline
 0 & 0 & 0 & 0 \\
\hline
 1+x^2+x y+x^2 y+z+y^2 z & 0 & 0 & x^2 y+z+y z+y^2 z \\
 x+y+x y+y^2 & 0 & 0 & 1+y+x y+y^2 \\
 z+y^2 z & 0 & 0 & x+y+y z+y^2 z \\
\hline
 0 & 1 & 0 & 0 \\
\hline
 y+x y+x y^2+y^2 z & 0 & 0 & 1+x+y+y^2+x y^2+z+y z+y^2 z
\end{array}\right),
\]
\[
 \sigma_{B4Z} =
\left(\begin{array}{cc|c|c}
 0 & 0 & 0 & 0 \\
\hline
 1+\frac{1}{y^2} & 1+\frac{1}{x} & 0 & 1+\frac{1}{x}+\frac{1}{y}+\frac{1}{y z} \\
\hline
 0 & 0 & 1 & 0 \\
\hline
 1+\frac{1}{y} & 1+\frac{1}{y} & 0 & 1+\frac{1}{z} \\
 1+\frac{1}{x}+\frac{1}{x y} & \frac{1}{z} & 0 & \frac{1}{x z}+\frac{1}{y z} \\
 \frac{1}{y^2} & \frac{1}{y} & 0 & 1+\frac{1}{x}+\frac{1}{z}+\frac{1}{y z} \\
\hline
 0 & 0 & 0 & 0 \\
\hline
 1+\frac{1}{y^2} & 0 & 0 & \frac{1}{x}+\frac{1}{y}+\frac{1}{z}+\frac{1}{y z}
\end{array}\right),
\]
\[
 r_{5X} =
\left(\begin{array}{cccccccc}
 1 & 0 & 0 & 0 & 0 & 0 & 0 & 0 \\
 1+x+y+x y+z & 1 & 0 & 0 & 0 & 0 & 1+y & 0 \\
 1+x+y & y & 1 & 1 & 1+x & 1+y & x+y & 1+y \\
 1+x^2+y z & 0 & 0 & 1 & 0 & 0 & 1+x & 0 \\
 x+y & 0 & 0 & 0 & 1 & 0 & 0 & 0 \\
 z+y z & 0 & 0 & 0 & 0 & 1 & 0 & 0 \\
 0 & 0 & 0 & 0 & 0 & 0 & 1 & 0 \\
 y+x y+z+y z & 0 & 0 & 0 & 0 & 0 & y & 1
\end{array}\right),
\]
\[
 r_{5Z} =
\left(\begin{array}{cccccccc}
 1 & 1+\frac{1}{x}+\frac{1}{y}+\frac{1}{x y}+\frac{1}{z} & \frac{1}{x y} & 1+\frac{1}{x^2}+\frac{1}{y z} & \frac{1}{x}+\frac{1}{y} & \frac{1}{z}+\frac{1}{y z} & 0 & \frac{1}{y}+\frac{1}{x y}+\frac{1}{z}+\frac{1}{y z} \\
 0 & 1 & \frac{1}{y} & 0 & 0 & 0 & 0 & 0 \\
 0 & 0 & 1 & 0 & 0 & 0 & 0 & 0 \\
 0 & 0 & 1 & 1 & 0 & 0 & 0 & 0 \\
 0 & 0 & 1+\frac{1}{x} & 0 & 1 & 0 & 0 & 0 \\
 0 & 0 & 1+\frac{1}{y} & 0 & 0 & 1 & 0 & 0 \\
 0 & 1+\frac{1}{y} & 1+\frac{1}{y} & 1+\frac{1}{x} & 0 & 0 & 1 & \frac{1}{y} \\
 0 & 0 & 1+\frac{1}{y} & 0 & 0 & 0 & 0 & 1
\end{array}\right),
\]
\[
 c_{5X} = 
\left(\begin{array}{cccc}
 y & 0 & 1 & 1+y \\
 1+x+x y & 1 & 1+x & y+x y \\
 1+y & 0 & 1 & y \\
 0 & 0 & 0 & 1
\end{array}\right),
\quad
c_{5Z} =
\left(\begin{array}{cccc}
 \frac{1}{y} & 0 & 1 & \frac{1}{z} \\
 0 & 1 & 0 & 1 \\
 1+\frac{1}{y} & 0 & 1 & \frac{1}{z} \\
 0 & 0 & 0 & 1
\end{array}\right).
\]
Factor out the first, third, and seventh qubits from $\sigma_{B4}$,
and coarse-grain by blocking two sites along $y$-direction.
\[
\sigma_{B4} \xrightarrow{y^2 \to y} \quad 
\sigma_{B5} = \begin{pmatrix} \sigma_{B5X} & 0 \\ 0 & \sigma_{B5Z} \end{pmatrix}, \text{ where }
\]
\[
\sigma_{B5X} =
\left(\begin{array}{cccc}
 1+x+x y+z & 1+z & y+x y & x+z \\
 y+y z & 1+x+x y+z & x y+y z & y+x y \\
 1+x^2+z+y z & x+x^2 & z+y z & x^2+z \\
 x y+x^2 y & 1+x^2+z+y z & x^2 y+y z & z+y z \\
 x+y & 1+x & 1+y & 1+x \\
 y+x y & x+y & y+x y & 1+y \\
 z+y z & 0 & x+y z & 1+z \\
 0 & z+y z & y+y z & x+y z \\
% x y+y z & 1+x & 1+x+y+x y+z+y z & 1+z \\
% y+x y & x y+y z & y+y z & 1+x+y+x y+z+y z
 x y+y z & 1+x & (1+x+z)(1+y) & 1+z \\
 y+x y & x y+y z & y+y z & (1+x+z)(1+y)
\end{array}\right),
\]
\[
\sigma_{B5Z} =
\left(\begin{array}{cccccc}
 1+\frac{1}{y} & 0 & 1+\frac{1}{x} & 0 & 1+\frac{1}{x} & \frac{1}{y}+\frac{1}{y z} \\
 0 & 1+\frac{1}{y} & 0 & 1+\frac{1}{x} & 1+\frac{1}{z} & 1+\frac{1}{x} \\
 1 & \frac{1}{y} & 1 & \frac{1}{y} & 1+\frac{1}{z} & 0 \\
 1 & 1 & 1 & 1 & 0 & 1+\frac{1}{z} \\
 1+\frac{1}{x} & \frac{1}{x y} & \frac{1}{z} & 0 & \frac{1}{x z} & \frac{1}{y z} \\
 \frac{1}{x} & 1+\frac{1}{x} & 0 & \frac{1}{z} & \frac{1}{z} & \frac{1}{x z} \\
 \frac{1}{y} & 0 & 0 & \frac{1}{y} & 1+\frac{1}{x}+\frac{1}{z} & \frac{1}{y z} \\
 0 & \frac{1}{y} & 1 & 0 & \frac{1}{z} & 1+\frac{1}{x}+\frac{1}{z} \\
 1+\frac{1}{y} & 0 & 0 & 0 & \frac{1}{x}+\frac{1}{z} & \frac{1}{y}+\frac{1}{y z} \\
 0 & 1+\frac{1}{y} & 0 & 0 & 1+\frac{1}{z} & \frac{1}{x}+\frac{1}{z}
\end{array}\right).
\]
Apply elementary symplectic transformations to factor out two more qubits.
\[
\sigma_{B6} =
 \begin{pmatrix} \sigma_{B6X} & 0 \\ 0 & \sigma_{B6Z} \end{pmatrix}
=
\begin{pmatrix} r_{6X} & 0 \\ 0 & r_{6Z} \end{pmatrix}
\begin{pmatrix} \sigma_{B5X} & 0 \\ 0 & \sigma_{B5Z} \end{pmatrix}
\begin{pmatrix} \id_{4 \times 4} & 0 \\ 0 & c_{6Z} \end{pmatrix}
\]
where
\[
 \sigma_{B6X} =
\left(\begin{array}{cccc}
 1+x+x y+z & 1+z & y+x y & x+z \\
 y+y z & 1+x+x y+z & x y+y z & y+x y \\
\hline
 0 & 0 & 0 & 0 \\
\hline
 x y+x^2 y & 1+x^2+z+y z & x^2 y+y z & z+y z \\
 x+y & 1+x & 1+y & 1+x \\
 y+x y & x+y & y+x y & 1+y \\
 z+y z & 0 & x+y z & 1+z \\
\hline
 0 & 0 & 0 & 0 \\
\hline
% x y+y z & 1+x & 1+x+y+x y+z+y z & 1+z \\
% y+x y & x y+y z & y+y z & 1+x+y+x y+z+y z % original formula from Mathematica
 x y+y z & 1+x & (1+x+z)(1+y) & 1+z \\
 y+x y & x y+y z & y+y z & (1+x+z)(1+y) % manual factoring
\end{array}\right),
\]
\[
 \sigma_{B6Z} =
\left(\begin{array}{c|c|c|ccc}
 0 & \frac{1}{y}+\frac{1}{x y} & 0 & \frac{1}{y^2}+\frac{1}{y} & \frac{1}{x}+\frac{1}{y}+\frac{1}{z}+\frac{1}{x z} & \frac{1}{x^2}+\frac{1}{x}+\frac{1}{x y}+\frac{1}{x z} \\
 0 & 1+\frac{1}{y} & 0 & 1+\frac{1}{x} & 1+\frac{1}{z} & 1+\frac{1}{x} \\
\hline
 1 & 0 & 0 & 0 & 0 & 0 \\
\hline
 0 & 1+\frac{1}{y} & 0 & 1+\frac{1}{y} & 1+\frac{1}{z} & 1+\frac{1}{z} \\
 0 & \frac{1}{x y}+\frac{1}{y z} & 0 & \frac{1}{y}+\frac{1}{x y} & 1+\frac{1}{x}+\frac{1}{z^2}+\frac{1}{x z} & 1+\frac{1}{x^2}+\frac{1}{z^2}+\frac{1}{y z} \\
 0 & 1+\frac{1}{x} & 0 & \frac{1}{x y}+\frac{1}{z} & \frac{1}{x}+\frac{1}{z} & \frac{1}{x^2}+\frac{1}{x} \\
 0 & 0 & 0 & \frac{1}{y^2}+\frac{1}{y} & 1+\frac{1}{x}+\frac{1}{y}+\frac{1}{z} & \frac{1}{y}+\frac{1}{x y} \\
\hline
 0 & 0 & 1 & 0 & 0 & 0 \\
\hline
 0 & 0 & 0 & \frac{1}{y^2}+\frac{1}{y} & 1+\frac{1}{x}+\frac{1}{y}+\frac{1}{z} & 1+\frac{1}{x}+\frac{1}{x y}+\frac{1}{z} \\
 0 & 1+\frac{1}{y} & 0 & 0 & 1+\frac{1}{z} & \frac{1}{x}+\frac{1}{z}
\end{array}\right),
\]
\[
 r_{6X} =
\left(\begin{array}{cccccccccc}
 1 & 0 & 0 & 0 & 0 & 0 & 0 & 0 & 0 & 0 \\
 0 & 1 & 0 & 0 & 0 & 0 & 0 & 0 & 0 & 0 \\
 1+y & 0 & 1 & 1 & 1+x & x & y & 0 & 1+y & 0 \\
 0 & 0 & 0 & 1 & 0 & 0 & 0 & 0 & 0 & 0 \\
 0 & 0 & 0 & 0 & 1 & 0 & 0 & 0 & 0 & 0 \\
 0 & 0 & 0 & 0 & 0 & 1 & 0 & 0 & 0 & 0 \\
 0 & 0 & 0 & 0 & 0 & 0 & 1 & 0 & 0 & 0 \\
 x+y & 0 & 0 & 0 & 1+x+z & x & y & 1 & 1+y & 0 \\
 0 & 0 & 0 & 0 & 0 & 0 & 0 & 0 & 1 & 0 \\
 0 & 0 & 0 & 0 & 0 & 0 & 0 & 0 & 0 & 1
\end{array}\right),
\]
\[
 r_{6Z} =
\left(\begin{array}{cccccccccc}
 1 & 0 & 1+\frac{1}{y} & 0 & 0 & 0 & 0 & \frac{1}{x}+\frac{1}{y} & 0 & 0 \\
 0 & 1 & 0 & 0 & 0 & 0 & 0 & 0 & 0 & 0 \\
 0 & 0 & 1 & 0 & 0 & 0 & 0 & 0 & 0 & 0 \\
 0 & 0 & 1 & 1 & 0 & 0 & 0 & 0 & 0 & 0 \\
 0 & 0 & 1+\frac{1}{x} & 0 & 1 & 0 & 0 & 1+\frac{1}{x}+\frac{1}{z} & 0 & 0 \\
 0 & 0 & \frac{1}{x} & 0 & 0 & 1 & 0 & \frac{1}{x} & 0 & 0 \\
 0 & 0 & \frac{1}{y} & 0 & 0 & 0 & 1 & \frac{1}{y} & 0 & 0 \\
 0 & 0 & 0 & 0 & 0 & 0 & 0 & 1 & 0 & 0 \\
 0 & 0 & 1+\frac{1}{y} & 0 & 0 & 0 & 0 & 1+\frac{1}{y} & 1 & 0 \\
 0 & 0 & 0 & 0 & 0 & 0 & 0 & 0 & 0 & 1
\end{array}\right),
\]
\[
 c_{6Z} =
\left(\begin{array}{cccccc}
 1 & 0 & 1 & \frac{1}{y} & 1 & 1+\frac{1}{x}+\frac{1}{z} \\
 0 & 1 & 0 & 0 & 0 & 0 \\
 0 & \frac{1}{y} & 1 & 0 & \frac{1}{z} & 1+\frac{1}{x}+\frac{1}{z} \\
 0 & 0 & 0 & 1 & 0 & 0 \\
 0 & 0 & 0 & 0 & 1 & 0 \\
 0 & 0 & 0 & 0 & 0 & 1
\end{array}\right).
\]
The third and eighth qubits in $\sigma_{B6}$ are trivial.
We finally finish our lengthy transformation.
\[
 \sigma_{B7} =
 \begin{pmatrix} \sigma_{B7X} & 0 \\ 0 & \sigma_{B7Z} \end{pmatrix}
=
\begin{pmatrix} r_{7X} & 0 \\ 0 & r_{7Z} \end{pmatrix}
\left( \begin{pmatrix} \sigma_{B6X} & 0 \\ 0 & \sigma_{B6Z} \end{pmatrix}\rvert_{\{3,8\}^c} \right)
\begin{pmatrix} c_{7X} & 0 \\ 0 & c_{7Z} \end{pmatrix}
\]
where
\[
 \sigma_{B7X} =
\left(\begin{array}{cccc}
 x+z & 1+x & 0 & 0 \\
 1+x & 1+z & 0 & 0 \\
 x+y & 1+y & 0 & 0 \\
 1+y & 1+x & 0 & 0 \\
 0 & 0 & x+z & 1+x \\
 0 & 0 & 1+x & 1+z \\
 0 & 0 & x+y & 1+y \\
 0 & 0 & 1+y & 1+x
\end{array}\right), \quad
\sigma_{B7Z} =
\left(\begin{array}{cccc}
 1+\frac{1}{y} & \frac{1}{x}+\frac{1}{y} & 0 & 0 \\
 1+\frac{1}{x} & 1+\frac{1}{y} & 0 & 0 \\
 1+\frac{1}{x} & \frac{1}{x}+\frac{1}{z} & 0 & 0 \\
 1+\frac{1}{z} & 1+\frac{1}{x} & 0 & 0 \\
 0 & 0 & 1+\frac{1}{y} & \frac{1}{x}+\frac{1}{y} \\
 0 & 0 & 1+\frac{1}{x} & 1+\frac{1}{y} \\
 0 & 0 & 1+\frac{1}{x} & \frac{1}{x}+\frac{1}{z} \\
 0 & 0 & 1+\frac{1}{z} & 1+\frac{1}{x}
\end{array}\right),
\]
\[
 r_{7X} =
\left(\begin{array}{cccccccc}
 1+z & 1 & 0 & 0 & z & 0 & 0 & 1 \\
 1+x & 1 & 1 & x & 0 & 1 & 0 & 1 \\
 1+x+y & 0 & 1 & x & 1 & 1+y & y & 1 \\
 1+y & 1 & 1 & 0 & 1+x & y & y & 1 \\
 1+x & 1 & 1 & x+z & 0 & 0 & 0 & 1 \\
 x & 1 & 1 & x+z & 0 & 1 & 1 & 1 \\
 1+x & 1 & 1 & 1+x+z & 1 & 1 & 1 & 1 \\
 0 & 0 & 0 & 1 & 0 & 0 & 0 & 0
\end{array}\right),
\]
\[
r_{7Z} =
\left(\begin{array}{cccccccc}
 0 & 0 & 1 & 0 & 0 & 0 & 0 & 1 \\
 0 & 1 & 0 & 0 & 0 & 1 & 1 & 1 \\
 0 & 1 & 0 & 0 & 0 & 0 & 0 & 1 \\
 1 & 1 & \frac{1}{x} & 0 & 1 & 0 & 1 & 0 \\
 1+\frac{1}{y} & 1 & 1+\frac{1}{x}+\frac{1}{x y} & 0 & 1+\frac{1}{y} & 1 & 1+\frac{1}{y} & \frac{1}{y} \\
 1+\frac{1}{x}+\frac{1}{y} & \frac{1}{x} & \frac{1}{x^2}+\frac{1}{x}+\frac{1}{x y}+\frac{1}{z} & 0 & \frac{1}{x}+\frac{1}{y} & 0 & \frac{1}{x}+\frac{1}{y} & 1+\frac{1}{y}+\frac{1}{z} \\
 1+\frac{1}{x} & \frac{1}{x} & \frac{1}{x^2}+\frac{1}{x}+\frac{1}{z} & 0 & \frac{1}{x} & 0 & 1+\frac{1}{x} & 1+\frac{1}{z} \\
 1+\frac{1}{z} & \frac{1}{z} & \frac{1}{x z} & 1 & \frac{1}{z} & \frac{1}{z} & 1 & 1+\frac{1}{z}
\end{array}\right),
\]
\[
 c_{7X} =
\left(\begin{array}{cccc}
 1 & 0 & 0 & 1+z \\
 0 & 1 & 0 & 1 \\
 1 & 0 & 1 & 1+z \\
 1 & 1 & 0 & 1+z
\end{array}\right),\quad
c_{7Z} =
\left(\begin{array}{cccc}
 \frac{1}{z} & 1+\frac{1}{z} & 1 & 0 \\
 1 & 1 & 0 & 0 \\
 \frac{1}{y} & 1+\frac{1}{y} & 0 & 1 \\
 0 & 1 & 0 & 0
\end{array}\right).
\]
We see that $\sigma_{B7}$ is a direct sum of two copies of $\sigma_B$.
This verifies the bifurcation:
\[
 A \xrightarrow{2\times 2 \times 2} A \oplus B, \text{ and } B \xrightarrow{2 \times 2 \times 2} B \oplus B.
\]

\section{Thermal partition function of the cubic code}
\label{sec:thermal-partition-function}

The only relevant property of the cubic code in relation to a thermal partition function
is that the generating map $\sigma$ for the stabilizer module is injective.
\[
 0 \to G \xrightarrow{\sigma} P
\]
This property is shared with 1D Ising model (Example~\ref{eg:1d-ising}),
2D toric code (Example~\ref{eg:2d-toric}), and 3D Chamon model (Example~\ref{eg:ChamonModel}).
Slightly more generally, Lemma~\ref{lem:coker-epsilon-resolution-length-D} says that
for two-dimensional exact code Hamiltonians one can choose an injective generating map 
that gives rise to an equivalent Hamiltonian in the sense of Definition~\ref{defn:equiv-H}.

Let $H = -\sum_i P_i$ be the Hamiltonian of the cubic code on $L \times L \times L$ periodic lattice.
We set the coupling constant to be $1$ so that each term literally squares to the identity $P_i^2 = 1$.
The thermal partition function is
\[
 \calZ = \calZ(\beta) = \trace \exp ( - \beta H )
\]
where $\beta$ is an inverse temperature.
Since $H$ consists of commuting terms,
the exponential function can be written as products.
\[
 \calZ = \trace \prod_i (\cosh \beta + P_i \sinh \beta) 
 = \cosh^M \beta \trace \prod_i (1 + x P_i)
\]
where $M=tL^3$ is the number of terms in $H$ and $x = \tanh \beta$.
Since the trace of a Pauli operator that is not identity is zero,
only the term in the expansion of the product 
that is proportional to the identity contributes to $\calZ$.
If there are $N=qL^3$ qubits in the system, 
the trace of the identity operator is $2^N$.
\[
 \frac{\calZ}{2^N \cosh^M \beta} = \sum_{\gamma \in \Gamma} x^{|\gamma|}
\]
where $\gamma$ runs over all possible collections of terms $P_i$ of $H$
that multiply to the identity. $\Gamma$ contains the empty collection.
Thus, $\Gamma$ is in one-to-one correspondence with $\ker \sigma_L$.
(It may be nonzero due to the periodic boundary conditions
although $\sigma$ was injective. See Section~\ref{sec:degeneracy}.)
Therefore, the cardinality of $\Gamma$ is precisely
\[
 |\Gamma| = 2^{\dim_{\FF_2} \ker \sigma_L} = 2^{k} ,
\]
the ground-state degeneracy.
The second equality is implied by Corollary~\ref{cor:k-formulas}.

We know from Corollary~\ref{cor:cubic-code-degeneracy-formula} that $k$ is bounded by
a linear function of $L$ for the cubic code; $k = O(L)$.
Therefore, the partition function is sandwiched as
\[
 | \log \calZ(\beta) - L^3\log (2^q \cosh^t \beta) | \le k \log 2 = O(L).
\]
Thus, the free energy per unit volume in the thermodynamic limit
is just a smooth function $\log( 2^q \cosh^t \beta )$ for all nonzero temperatures.
This should contrast with three- or higher-dimensional toric code model 
where the free energy density has a singularity at a nonzero temperature~\cite{NussinovOrtiz2008Autocorrelations}.

The analyticity of the free energy does not directly invalidate
a possibility of self-correction of encoded quantum information.
The latter is rather a dynamical process,
which may have little to do with the thermal equilibrium.
In fact, it is quite subtle to analyze self-correcting power
or to give a criterion on it, based on the thermal partition function.
The existence or possibility of a good decoding process
to extract the encoded quantum information
would be a more appropriate way to address the question of self-correction.
For instance, in the four-dimensional toric code model
the thermal expectation value for any bare logical operator is zero.
However, a special dressed logical operator defined with respect to a decoding algorithm
can have a non-vanishing expectation value in the thermodynamic
limit~\cite{ChesiLossBravyiEtAl2010Thermodynamic}.
For the cubic code, the analyticity of the partition function suggests that
it does not allow such a nonzero thermal expectation value for logical operators.
Nevertheless, we can show that the characteristic time scale
of the expectation values of logical operators dressed by a decoding algorithm,
is very large at low temperatures.
This is the topic of the next chapters.